% Part: first-order-logic
% Chapter: axiomatic-deduction
% Section: provability-consistency

\documentclass[../../../include/open-logic-section]{subfiles}

\begin{document}

\iftag{FOL}
      {\olfileid{fol}{axd}{prv}}
      {\olfileid{pl}{axd}{prv}}

\olsection{\usetoken{S}{derivability} आणि सुसंगतता}

आता आपण !!{derivability} संबंधाचे अनेक गुणधर्म सिद्ध करू. ते स्वतंत्रपणेही
महत्त्वाचे आहेत, पण संपूर्णता प्रमेयाच्या सिद्धतेत प्रत्येकाची भूमिका असेल.

\begin{prop}\ollabel{prop:provability-contr}
  जर $\Gamma \Proves !A$ आणि $\Gamma \cup \{!A\}$ विसंगत असेल, तर
  $\Gamma$ विसंगत आहे.
\end{prop}

\begin{proof}
  जर $\Gamma \cup \{!A\}$ विसंगत असेल, तर $\Gamma \cup \{!A\}
  \Proves \lfalse$. \olref[ptn]{prop:reflexivity} वरून प्रत्येक
  $!B \in \Gamma$ साठी $\Gamma \Proves !B$. गृहीतकानुसार
  $\Gamma \Proves !A$ सुद्धा असल्यामुळे प्रत्येक $!B \in \Gamma \cup
  \{!A\}$ साठी $\Gamma \Proves !B$. \olref[ptn]{prop:transitivity}
  वरून $\Gamma \Proves \lfalse$, म्हणजे $\Gamma$ विसंगत आहे.
\end{proof}

\begin{prop}
\ollabel{prop:prov-incons}
$\Gamma \Proves !A$ तेव्हा आणि केवळ तेव्हाच, जेव्हा $\Gamma \cup \{\lnot !A\}$
विसंगत असते.
\end{prop}

\begin{proof}
प्रथम $\Gamma \Proves !A$ असे समजा. मग \olref[ptn]{prop:monotonicity}
वरून $\Gamma \cup \{\lnot !A\} \Proves !A$. \olref[ptn]{prop:reflexivity}
वरून $\Gamma \cup \{\lnot !A\} \Proves \lnot !A$. तसेच
\olref[prp]{ax:lnot2} वरून $\Proves \lnot !A \lif (!A \lif \lfalse)$.
म्हणून \olref[ded]{prop:mp} चे दोन उपयोग करून $\Gamma \cup \{\lnot
!A\} \Proves \lfalse$ मिळते.

आता $\Gamma \cup \{\lnot !A\}$ विसंगत आहे असे गृहीत धरा, म्हणजे
$\Gamma \cup \{\lnot !A\} \Proves \lfalse$. निगमन प्रमेयावरून $\Gamma
\Proves \lnot !A \lif \lfalse$. \olref[prp]{ax:lfalse2} वरून $\Gamma
\Proves (\lnot !A \lif \lfalse) \lif \lnot\lnot !A$, म्हणून
\olref[ded]{prop:mp} वरून $\Gamma \Proves \lnot\lnot !A$. आपल्याकडे
$\Gamma \Proves \lnot\lnot !A \lif !A$ हेही असल्यामुळे
(\olref[prp]{ax:dne}), \olref[ded]{prop:mp} च्या आणखी एका उपयोगाने
$\Gamma \Proves !A$ मिळते.
\end{proof}

\begin{prob}
$\Gamma \Proves \lnot !A$ तेव्हा आणि केवळ तेव्हाच, जेव्हा $\Gamma \cup \{!A\}$
विसंगत असते, हे सिद्ध करा.
\end{prob}

\begin{prop}\ollabel{prop:explicit-inc}
  जर $\Gamma \Proves !A$ आणि $\lnot !A \in \Gamma$ असेल, तर $\Gamma$
  विसंगत आहे.
\end{prop}

\begin{proof}
  \olref[prp]{ax:lnot2} वरून $\Gamma \Proves \lnot !A \lif (!A \lif \lfalse)$.
  \olref[ded]{prop:mp} चे दोन उपयोग करून $\Gamma \Proves \lfalse$ मिळते.
\end{proof}


\begin{prop}\ollabel{prop:provability-exhaustive}
  $\Gamma \cup \{!A\}$ आणि $\Gamma \cup \{\lnot !A\}$ हे दोन्ही संच
  विसंगत असतील, तर $\Gamma$ विसंगत आहे.
\end{prop}

\begin{proof}
सराव.
\end{proof}

\tagprob{FOL}
\begin{prob}
  \olref[fol][axd][prv]{prop:provability-exhaustive} सिद्ध करा
\end{prob}
\tagendprob

\tagprob{notFOL}
\begin{prob}
  \olref[pl][axd][prv]{prop:provability-exhaustive} सिद्ध करा
\end{prob}
\tagendprob


\end{document}
